\chapter{KAJIAN PUSTAKA}

\section{Manajemen Tugas dan Waktu Akademik}
Manajemen waktu adalah proses merencanakan dan mengendalikan penggunaan
waktu untuk mencapai tujuan tertentu. Dalam konteks akademik, perilaku
manajemen waktu meliputi penetapan tujuan dan prioritas, mekanisme
perencanaan (membuat daftar tugas dan jadwal), serta preferensi terhadap
keteraturan. Penelitian pada mahasiswa di Indonesia menunjukkan bahwa
perilaku manajemen waktu berpengaruh positif dan signifikan terhadap
prestasi akademik mahasiswa \parencite{kamilatunnisa2025}.

Di sisi lain, prokrastinasi akademik --- menunda pekerjaan meskipun sadar
akan konsekuensinya --- merupakan kegagalan regulasi diri yang umum
terjadi pada mahasiswa; semakin rendah regulasi diri mahasiswa, semakin
tinggi kecenderungan prokrastinasi akademiknya \parencite{fadila2024}. Salah
satu pemicunya adalah beban kognitif dalam
mengingat dan mengorganisasi banyak tugas dari berbagai sumber. Sistem
yang mampu mengumpulkan, memecah, dan menjadwalkan tugas secara otomatis
berpotensi menurunkan beban tersebut dan mendorong perilaku belajar yang
lebih teratur.

\section{Kecerdasan Buatan dan Model Bahasa Besar}
Kecerdasan buatan (\emph{artificial intelligence}) adalah bidang ilmu yang
mempelajari sistem cerdas yang mampu memahami lingkungannya dan mengambil
tindakan untuk mencapai tujuan tertentu. Salah satu inovasinya yang
berkembang paling pesat adalah model bahasa besar (\emph{large language
model}/LLM), yaitu model bahasa berbasis AI yang mampu memahami dan
menghasilkan teks dalam bahasa alami \parencite{johan2026}.

Kajian terhadap pemanfaatan LLM di perguruan tinggi menunjukkan bahwa LLM
dapat membantu mahasiswa memahami materi perkuliahan, memperoleh
informasi dengan cepat, mendukung pembelajaran mandiri, serta
meningkatkan interaksi dalam proses belajar apabila digunakan secara
bijak dan didukung literasi digital yang memadai \parencite{johan2026}.
Chatbot berbasis AI juga terbukti berperan sebagai asisten virtual yang
personal, responsif, dan adaptif dalam pembelajaran digital, termasuk
mendukung pengembangan keterampilan manajemen waktu dan kemandirian
belajar \parencite{hidayah2025}.
Dalam Welya, kemampuan ini dimanfaatkan untuk tiga hal: (1) menyaring
email kampus yang mengandung penugasan, (2) mengekstraksi atribut tugas
seperti judul, mata kuliah, dan tenggat, serta (3) memecah tugas menjadi
subtask dengan estimasi waktu pengerjaan.

\section{Aplikasi Web, Aplikasi Desktop, dan Integrasi Layanan Google}
Aplikasi web dipilih karena dapat diakses lintas perangkat tanpa
instalasi, sementara versi desktop dibangun menggunakan kerangka kerja
Tauri yang membungkus antarmuka web yang sama menjadi aplikasi native
yang ringan. Penelitian sejenis menunjukkan bahwa aplikasi manajemen tugas
berbasis web dengan fitur pengingat otomatis, prioritas, dan kalender
interaktif mampu meningkatkan disiplin dan efisiensi pengelolaan waktu
mahasiswa \parencite{panjaitan2025}.

Integrasi dengan ekosistem Google dilakukan melalui API resmi Gmail,
Google Tasks, dan Google Calendar. Pemanfaatan Google Calendar untuk
otomatisasi penjadwalan dan pengingat tenggat terbukti meningkatkan
kepatuhan terhadap tenggat waktu sekaligus mengurangi kesalahan
penjadwalan manual \parencite{setiawan2025}, dan integrasi Google Calendar
API mampu menyinkronkan tugas ke kalender pribadi pengguna secara
\emph{real-time} \parencite{moedjahedy2025}. Otorisasi akses data pengguna
menggunakan protokol OAuth 2.0, standar industri yang memungkinkan
aplikasi pihak ketiga mengakses sumber daya pengguna secara terbatas
tanpa menyimpan kata sandi; penerapannya terbukti meningkatkan keamanan
autentikasi secara signifikan \parencite{arianto2025}. Dengan pendekatan ini,
data
tetap berada di akun Google milik pengguna dan aplikasi hanya meminta
cakupan izin yang diperlukan.

Karena sistem memproses isi email pengguna, aspek privasi menjadi
perhatian utama dalam perancangan. Akses data hanya dilakukan setelah
persetujuan eksplisit pengguna, cakupan izin (\emph{scope}) OAuth
dibatasi seminimal mungkin, data yang disimpan hanya hasil ekstraksi
tugas --- bukan salinan isi email --- dan pengguna dapat mencabut akses
kapan saja. Praktik ini sejalan dengan prinsip pelindungan data pribadi
sebagaimana diatur dalam Undang-Undang Nomor 27 Tahun 2022 tentang
Pelindungan Data Pribadi.

\section{\emph{System Usability Scale} (SUS)}
SUS adalah instrumen baku berisi sepuluh pernyataan berskala Likert untuk
mengukur kebergunaan suatu sistem secara cepat dan andal, dan telah
banyak digunakan dalam evaluasi aplikasi di Indonesia \parencite{dyayu2023}.
Skor akhir berada pada rentang 0--100 dengan nilai 68 sebagai ambang
rata-rata industri. Instrumen ini digunakan pada tahap evaluasi untuk
mengukur kemudahan penggunaan Welya oleh mahasiswa.

\section{Penelitian dan Produk Sejenis}
Beberapa produk manajemen tugas telah tersedia secara luas, di antaranya
Google Tasks, Todoist, dan Notion. Ketiganya andal sebagai pencatat tugas,
tetapi memiliki keterbatasan dalam konteks akademik mahasiswa Indonesia,
sebagaimana diringkas pada Tabel~\ref{tab:perbandingan}.

\begin{table}[htbp]
  \centering
  \caption{Perbandingan Welya dengan produk sejenis}
  \label{tab:perbandingan}
  \small
  \begin{tabular}{@{}p{4.2cm}cccc@{}}
    \toprule
    Fitur & Google Tasks & Todoist & Notion & Welya \\
    \midrule
    Pencatatan tugas manual        & $\checkmark$ & $\checkmark$ & $\checkmark$ & $\checkmark$ \\
    Deteksi tugas dari email       & --           & --           & --           & $\checkmark$ \\
    Subtask + estimasi waktu (AI)  & --           & --           & --           & $\checkmark$ \\
    Kalender jadwal kuliah         & --           & $\checkmark$ & $\checkmark$ & $\checkmark$ \\
    Pencarian cerdas (AI)          & --           & --           & $\checkmark$ & $\checkmark$ \\
    Berfokus konteks akademik      & --           & --           & --           & $\checkmark$ \\
    \bottomrule
  \end{tabular}
\end{table}

Kesenjangan (\emph{gap}) yang mendasari inovasi ini adalah belum adanya
asisten akademik yang mendeteksi tugas secara otomatis dari email kampus
berbahasa Indonesia, memecahnya menjadi subtask beserta estimasi waktu,
dan menyatukannya dengan jadwal kuliah dalam satu dasbor. Welya mengisi
kesenjangan tersebut dengan menggabungkan LLM dan integrasi layanan
Google dalam satu produk yang dirancang khusus untuk alur kerja
mahasiswa.
